\chapter{Cost and latency}
\label{ch:cost}

\section*{What this chapter is for}

\reported{Cost and latency are among the most frequently asked topics in this
domain.}{field-guide} They are also the surface where a senior backend
engineer has the largest unfair advantage, because the reasoning is capacity
planning and almost nobody arriving from a data-science route has done it.

\section*{The questions}

\begin{bankq}{What does this cost to run?}
\qasked
Usually attached to a system just designed on a whiteboard, and often phrased
casually.

\qtested
Whether you can do arithmetic out loud without a spreadsheet, and whether you
know which quantities to ask for.

\qstaff
Ask for the three numbers, then multiply in front of them: how many users, how
many interactions each, how much text per interaction. From those, tokens per
day; from tokens per day and a unit price, a daily figure; from a daily figure,
a monthly one \dash{} and the comparison that matters is against the revenue or
the budget of the thing it is part of.

\reported{The arithmetic is the expected form of the answer.}{field-guide-get-hired}

Then the parts that separate levels. Input and output are priced differently
and usually by a wide margin, so a system that reads a lot and writes a little
has a different shape from one that does the reverse, and the retrieval
decision in Chapter~\ref{ch:retrieval} moves this number more than the model
choice does. The conversation history is re-sent every turn unless you do
something about it, so a long conversation is quadratic in a way people do not
expect. And a retry is a full second payment.

\judgement{The sentence worth having ready: ``the number I would want first is
average context length, because it multiplies everything else and it is the one
nobody measures.''}

\qsenior
Naming the cost drivers correctly \dash{} model choice, context length,
caching \dash{} without producing a figure. It is right and it does not answer
the question, which was a number.

\qcontested
Unit prices move constantly and any figure in a book is wrong. \judgement{Which
is precisely why the method is the answer: learn to ask for the three
quantities and multiply, and the arithmetic survives every price change.}

\qdrill
Do it now, out loud, for a system you know: users, interactions, tokens, price,
day, month. Under two minutes. It is much harder the first time than you expect
and much easier the second, which is the whole reason to do it once before
somebody asks.
\end{bankq}

\begin{bankq}{It is too slow. What do you do?}
\qasked
A scenario, usually with a target attached.

\qtested
Whether you have an order of attack or a list of ideas. Chapter~\ref{ch:answering}
step one applies harder here than anywhere: which latency, for whom.

\qstaff
Establish what is being measured first. Time to first token and time to
completion are different numbers with different fixes, and a streaming
interface makes the first one most of what users perceive. Average and tail are
different problems; a good average with a bad tail is the usual shape and the
tail is what people remember.

Then attack in order of what is free. Is the work necessary \dash{} can the
answer be cached, or does the request need a model at all? Is it parallel
\dash{} are there sequential calls that do not depend on each other? Is the
output shorter than it is \dash{} generation time is roughly linear in output
length and truncating the response is often the cheapest real win. Is a smaller
model sufficient for this step \dash{} not for the product, for this step. And
only then, is the prompt smaller.

\judgement{The structural answer that reads as senior: the largest latency wins
in these systems usually come from removing a call rather than from speeding
one up. A pipeline of four sequential model calls has four times the tail risk
of one, and the design question is whether it needs to be four.}

\qsenior
Caching, a smaller model, a shorter prompt, streaming. All correct, offered in
parallel, with no measurement first.

\qcontested
Whether to route between models by difficulty is a real design fork with real
adherents on both sides. \reported{It is reported as less prominent than the
other cost topics and not a dedicated interview subject.}{field-guide-system-design}
\judgement{Routing adds a classification step that has its own latency and its
own error rate, and the error rate is what decides whether it pays.}

\qdrill
For a system you know, name the four latency numbers \dash{} first token and
completion, average and tail \dash{} and say which one a user would complain
about. If you cannot, that is what to instrument.
\end{bankq}

\begin{bankq}{Where would caching help, and where is it a trap?}
\qasked
Often follows the latency question.

\qtested
Whether you apply caching reasoning you already have, or treat this as a new
subject.

\qstaff
It is the caching reasoning you already have, with one addition. The familiar
part: a cache pays when requests repeat, and a long tail of unique requests
gets you an eviction cost and no hits. Measure the repeat rate before building
anything.

The addition is that there are two distinct things to cache. Whole responses,
keyed on the request \dash{} which works when questions genuinely repeat, and
which has the same invalidation problem as any cache, with the wrinkle that a
stale answer here is not obviously stale to the user. And the repeated
\emph{prefix} of a request \dash{} the instructions, the retrieved documents,
the conversation so far \dash{} which several providers support directly and
which is nearly free to adopt because it does not change what the system
returns.

\judgement{The second is underused relative to how well it pays, because it
does not feel like an architectural decision. For a system with a long fixed
instruction and short user turns, it is frequently the single largest cost
reduction available and it costs an afternoon.}

The trap worth naming: caching a non-deterministic system freezes one sample of
its behaviour. If the cached answer was bad, it is now consistently bad, and a
user who tries again gets the same thing \dash{} which removes the one recovery
mechanism users have.

\qsenior
Cache repeated queries with a sensible expiry and be careful about
invalidation. Correct, and misses prefix caching entirely.

\qcontested
Semantic caching \dash{} serving a cached answer for a similar rather than
identical question \dash{} is genuinely disputed. \judgement{It trades a
correctness risk for a cost saving, and the risk is that similarity is not
equivalence. Reasonable for suggestions, unreasonable for anything a person
will act on.}

\qdrill
For a system you know, estimate what fraction of its request tokens are
identical across requests. If it is most of them, you have just found the
answer to the cost question two pages back.
\end{bankq}

\section*{What this chapter does not claim}

No prices, no benchmarks, no figures. They date within a quarter, and a book
carrying them would be actively misleading by the time it was read \dash{}
which is why the arithmetic is the content.

The ordering of latency attacks, the position on routing, and the claim about
prefix caching being underused are the author's judgement. The last is the one
most likely to become dated, as provider support and defaults change.
